\section{Introduction}
\label{sec:intro}

A taxonomy is a hierarchy of classes that organises a domain's vocabulary so
that machines and people can classify, browse and retrieve its objects
consistently. On the Web, taxonomies drive product catalogues, search
verticals and content moderation, and the community has studied how to expand
and complete them from Web text for years~\cite{shen2020taxoexpan,
wang2021hef,jiang2022taxoenrich,arous2023taxocomplete,shi2024icon}. Crypto-assets
are a domain in which the vocabulary itself is contested. Tokens are issued by
smart contracts, traded on exchanges and decentralised venues, and described in
whitepapers, explorers and social media, and the documents that say what a
token \emph{is} are themselves Web publications, among them the EU's Markets
in Crypto-Assets regulation (MiCA), the US GENIUS and CLARITY Acts, the UK's
cryptoasset regulations, and the guidance of the Financial Action Task Force
(FATF), the Basel Committee, the FSB and IOSCO, alongside academic
taxonomies~\cite{tasca2019taxonomy,zhang2026cryptotaxonomy}.

\paragraph{Relevance to the Web.}
This work falls under the Semantics and Knowledge track. The research challenge
is Web-specific because the definitions that govern crypto-assets are
published, versioned and consumed as Web documents by heterogeneous bodies,
no shared vocabulary links them, and the assets they describe exist only as records on
Web-accessible ledgers. The result is a machine-readable, provenance-tracked
vocabulary that block explorers, exchanges and compliance services can use to
classify tokens consistently across jurisdictions and to regenerate the
classification when an instrument changes.

\begin{figure}[t]
  \centering
  %% Figure 1: one real term, four sources, one induced term type (column width).
\begin{tikzpicture}[
  font=\footnotesize,
  src/.style={draw=gray!70, fill=gray!8, rounded corners=2pt, text width=2.75cm, inner sep=2.5pt, align=left},
  tt/.style={draw=blue!60!black, fill=blue!6, rounded corners=2pt, text width=3.25cm, inner sep=3pt, align=left},
  val/.style={font=\scriptsize, inner sep=1pt},
  arr/.style={-{Latex[length=1.6mm]}, gray!80, thin},
  node distance=2mm and 5mm]
  \node[src] (mica) {\textbf{MiCA} Art.~3(6--7)\\[1pt]\scriptsize ``e-money token \ldots\ referencing the value of \emph{one official currency}''; ``asset-referenced token \ldots\ referencing \emph{another value or right} or a combination''};
  \node[src, below=of mica] (genius) {\textbf{GENIUS Act} \S2(22)\\[1pt]\scriptsize ``maintain \ldots\ a stable value relative to \ldots\ a \emph{fixed amount of monetary value}''};
  \node[src, below=of genius] (uk) {\textbf{UK SI 2026/102} art.~88G\\[1pt]\scriptsize ``a stable value in relation to \emph{a particular fiat currency}''};
  \node[src, below=of uk] (basel) {\textbf{Basel SCO60.6}\\[1pt]\scriptsize Group~1b: ``cryptoassets with effective \emph{stabilisation mechanisms}''; Group~2: ``stablecoins and unbacked cryptoassets''};
  \node[tt, right=of $(genius.east)!0.5!(uk.east)$] (tt) {\textbf{Term type: reference asset}\\[1pt]\scriptsize\emph{What does the asset reference to maintain a stable value?}\\[3pt]
    \scriptsize
    \textbullet~single official currency (EU)\\
    \textbullet~basket or other assets (EU)\\
    \textbullet~fiat currency (US, UK)\\
    \textbullet~commodity (Basel)\\
    \textbullet~other crypto-assets / none (Basel)};
  \draw[arr] (mica.east) -- (tt.west|-mica.east) -- ([yshift=4mm]tt.west);
  \draw[arr] (genius.east) -- ([yshift=1mm]tt.west);
  \draw[arr] (uk.east) -- ([yshift=-1mm]tt.west);
  \draw[arr] (basel.east) -- (tt.west|-basel.east) -- ([yshift=-4mm]tt.west);
  \node[font=\scriptsize\itshape, gray!90, above=1mm of mica.north west, anchor=south west] {term: \emph{stablecoin}};
\end{tikzpicture}

  \caption{One term, four sources, one induced term type. Each instrument
  defines \term{stablecoin} by what the asset references, but the admissible
  references differ by jurisdiction. \sys makes the dimension explicit, keeps
  one value per jurisdictional position, and cites the section each value comes
  from.}
  \label{fig:one}
\end{figure}

These sources disagree in a specific way. \Cref{fig:one} shows the term
\term{stablecoin}. \src{MiCA} splits it into e-money tokens, which reference one
official currency, and asset-referenced tokens, which reference anything else.
The \src{GENIUS Act} and the UK regulations admit only a fiat currency as the
reference. The Basel Committee admits a commodity and treats algorithmic
stabilisation as a different category. None of the texts states that ``what the
asset references'' is a classification dimension, yet each uses it to draw a
legal or prudential boundary. We call such a dimension a \emph{term type}, a
question along which definitions discriminate, together with the set of values
the sources admit. The reference asset is one of many; the issuer's
identifiability, the fungibility of the asset, the custody model and the
consensus mechanism are others, and they recur across instruments under
different names.

Existing approaches do not recover term types from such a corpus. First,
manually built taxonomies of crypto-assets and of blockchain
architectures~\cite{tasca2019taxonomy,zhang2026cryptotaxonomy} encode one
perspective each, a technical or an investor's one, and cannot be regenerated
when an instrument changes. Second, LLM methods for taxonomy
completion~\cite{shi2024icon}, hierarchical
classification~\cite{zhang2025teleclass} and tagging~\cite{wan2024tntllm,
li2023taggpt} start from a label space that already exists, and ontology
learning with LLMs~\cite{babaeigiglou2023llms4ol} is evaluated on general-domain
benchmarks where the gold hierarchy is known. Third, none of these methods ties
an induced class to the passage that licenses it, which in a legal domain is the
difference between a taxonomy a compliance team can audit and one it cannot.

In this paper we propose \sys, a framework that induces term types from the
definitions in a corpus of legal instruments and academic taxonomies and
attaches provenance to every node. \sys builds a definitions-first corpus, in
which each document is reduced to the sections that define or classify; it
extracts term types in one pass over the whole corpus held in a cached context,
so that the model can recognise the same dimension under different names across
sources; it requires a document index, an exact section label and a verbatim
quote for every term type and every value; and it verifies each quote against
the corpus, repairing wrong section labels and rejecting quotes it cannot find.
It then orders the catalogue. For each term type it asks for a cited probability
distribution over the more general term types it specialises, and consolidates
the weighted edges greedily into a directed acyclic graph (DAG), so that a term
type can sit under more than one parent. The output is an ontology of term types
with values, aliases, jurisdictions, recorded conflicts and weighted is-a edges,
every element of which cites its source.

\smallskip\noindent\textbf{Contributions.}
\begin{itemize}
  \item We formalise term type induction, in which a corpus of definitional
    text yields the classification dimensions it uses, with values and
    provenance, without a seed taxonomy or gold labels (\cref{sec:problem}).
  \item We propose \sys, which combines a definitions-first corpus, single-pass
    extraction over a cached whole-corpus context with mandatory citations, a
    citation verifier, and a per-node is-a query consolidated into a DAG with
    metrics for its later maintenance (\cref{sec:method}).
  \item We release \corpus, a definitions-first corpus of 24 documents across
    four jurisdictions, and the resulting ontology of 93 term types with 317
    values, 480 citations and 279 weighted is-a edges
    (\cref{sec:exp:rq1,sec:exp:rq2}).
  \item We evaluate faithfulness and coverage on \corpus: 97.9\% of term-type
    citations verify, every document contributes to at least one term type, 69
    term types record a cross-source conflict, and the hierarchy resolves
    \BBits\ bits per root-to-leaf descent under one root per category. \todo{Add the headline comparative
    number from RQ2--RQ4.} Code and data are at an anonymised
    repository.\footnote{\todo{anonymised URL}}
\end{itemize}
